
\documentclass{article}
\usepackage{colortbl}
\usepackage{makecell}
\usepackage{multirow}
\usepackage{supertabular}

\begin{document}

\newcounter{utterance}

\centering \large Interaction Transcript for game `cladder', experiment `full\_v1.5\_default', episode 7573 with selfrefine{-}Qwen3.5{-}27B.
\vspace{24pt}

{ \footnotesize  \setcounter{utterance}{1}
\setlength{\tabcolsep}{0pt}
\begin{supertabular}{c@{$\;$}|p{.15\linewidth}@{}p{.15\linewidth}p{.15\linewidth}p{.15\linewidth}p{.15\linewidth}p{.15\linewidth}}
    \# & \multicolumn{2}{c}{Player} && \multicolumn{2}{c}{Game Master} \\
    \hline

    \theutterance \stepcounter{utterance}  
    & & & \multicolumn{4}{p{0.6\linewidth}}{
        \cellcolor[rgb]{0.9,0.9,0.9}{
            \makecell[{{p{\linewidth}}}]{
                \texttt{\tiny{[P1$\langle$GM]}}
                \texttt{You are an expert in causal inference. The following question is not a typical commonsense query, but rather a meticulously designed question created by a professor specializing in causal inference, intended to assess the students' mastery of the course content.} \\
\texttt{Imagine a self{-}contained, hypothetical world with only the following conditions, and without any unmentioned factors or causal relationships: Pexu has a direct effect on hwax and rukz. Hwax has a direct effect on rukz. For those who are not pexu, the probability of rukz is 17\%. For those who are pexu, the probability of rukz is 80\%. For those who are pexu, would it be more likely to see rukz if the individual was not pexu?} \\
\texttt{Start your answer with 'Yes' or 'No', followed by additional reasoning or evidence to support your explanation.} \\
            }
        }
    }
     \\ \\

    \theutterance \stepcounter{utterance}  
    & \multicolumn{4}{p{0.6\linewidth}}{
        \cellcolor[rgb]{0.561,0.737,0.561}{
            \makecell[{{p{\linewidth}}}]{
                \texttt{\tiny{[P1$\rangle$GM]}}
                \texttt{No} \\
\\ 
\texttt{The question asks whether, for an individual who is currently **pexu**, the probability of observing **rukz** would be higher if that same individual were **not pexu**. This is a counterfactual query comparing the outcome under the current state (\$P(\textbackslash{}text\{rukz\} \textbackslash{}mid \textbackslash{}text\{pexu\}=1)\$) with the outcome under the hypothetical intervention (\$P(\textbackslash{}text\{rukz\} \textbackslash{}mid \textbackslash{}text\{do\}(\textbackslash{}text\{pexu\}=0))\$).} \\
\\ 
\texttt{Based on the conditions provided:} \\
\texttt{1.  The probability of **rukz** given **pexu** is true is **80\%**.} \\
\texttt{2.  The probability of **rukz** given **pexu** is false (not pexu) is **17\%**.} \\
\\ 
\texttt{The causal structure indicates that **pexu** has a direct effect on **rukz** (and also an indirect effect via **hwax**). The data explicitly states that the presence of **pexu** increases the likelihood of **rukz** from 17\% to 80\%. Therefore, if an individual who is currently **pexu** were hypothetically changed to **not pexu**, the probability of **rukz** would drop from 80\% to 17\%.} \\
\\ 
\texttt{Consequently, it would be **less** likely, not more likely, to see **rukz** if the individual were not **pexu**. The premise of the question ("would it be more likely") is false.} \\
            }
        }
    }
    & & \\ \\

    \theutterance \stepcounter{utterance}  
    & & & \multicolumn{2}{p{0.3\linewidth}}{
        \cellcolor[rgb]{0.9,0.9,0.9}{
            \makecell[{{p{\linewidth}}}]{
                \texttt{\tiny{[GM$|$GM]}}
                \texttt{No} \\
            }
        }
    }
    & & \\ \\

    \theutterance \stepcounter{utterance}  
    & & & \multicolumn{2}{p{0.3\linewidth}}{
        \cellcolor[rgb]{0.9,0.9,0.9}{
            \makecell[{{p{\linewidth}}}]{
                \texttt{\tiny{[GM$|$GM]}}
                \texttt{no} \\
            }
        }
    }
    & & \\ \\

    \theutterance \stepcounter{utterance}  
    & & & \multicolumn{2}{p{0.3\linewidth}}{
        \cellcolor[rgb]{0.9,0.9,0.9}{
            \makecell[{{p{\linewidth}}}]{
                \texttt{\tiny{[GM$|$GM]}}
                \texttt{game\_result = WIN} \\
            }
        }
    }
    & & \\ \\

\end{supertabular}
}

\end{document}
